# 📘 PART II — Non‑Linear Operator Layer + Thermodynamic Layer

`latex
%------------------------------------------------------------
\section{Non-Linear Operator Layer}

\subsection{Activation Operators}

Activation operators induce curvature in the tensor manifold. For a state $x$,
the activation-induced curvature term is denoted $\Phi(x)$.

\paragraph{GELU.}
\[
\Phi_{\text{GELU}}(x) =
x \cdot \frac{1}{2}\left(1 + \mathrm{erf}\left(\frac{x}{\sqrt{2}}\right)\right).
\]

\paragraph{SwiGLU.}
\[
\Phi{\text{SwiGLU}}(x) = (xW1) \odot \sigma(xW_2),
\]
where $W1, W2$ are learned matrices and $\sigma$ is the sigmoid function.

\paragraph{Softplus.}
\[
\Phi_{\text{Softplus}}(x) = \log(1 + e^x).
\]

\paragraph{Tanh.}
\[
\Phi_{\tanh}(x) = \tanh(x).
\]

These operators generate non-linear curvature contributions that directly affect
holonomy transport and thermodynamic gradients.

%------------------------------------------------------------
\subsection{Normalization Operators}

Normalization stabilizes curvature and prevents drift blowout.

\paragraph{LayerNorm.}
\[
\mathrm{LayerNorm}(x) = \frac{x - \mu(x)}{\sigma(x)}.
\]

\paragraph{RMSNorm.}
\[
\mathrm{RMSNorm}(x) = \frac{x}{\sqrt{\frac{1}{n}\sumi xi^2}}.
\]

\paragraph{GroupNorm.}
Group normalization partitions $x$ into groups $G_k$:
\[
\mathrm{GroupNorm}(x) = \frac{x - \mu(Gk)}{\sigma(Gk)}.
\]

%------------------------------------------------------------
\subsection{Residual Mixing Operators}

Residual mixing preserves holonomy invariants and stabilizes geometric evolution.

\paragraph{Standard Residual.}
\[
\mathrm{Residual}(x, y) = x + y.
\]

\paragraph{Curvature-Aware Residual.}
\[
\mathrm{Residual}_{\kappa}(x, y) = x + \kappa(y).
\]

This operator ensures curvature consistency across update steps.

%------------------------------------------------------------
\subsection{Operator Algebra}

Let $\mathcal{O}$ denote the set of all engine-layer operators.

\paragraph{Composition.}
\[
(O2 \circ O1)(x) = O2(O1(x)).
\]

\paragraph{Associativity.}
\[
O3 \circ (O2 \circ O1) = (O3 \circ O2) \circ O1.
\]

\paragraph{Non-Commutation.}
\[
\mathrm{Norm} \circ \mathrm{Act} \ne \mathrm{Act} \circ \mathrm{Norm}.
\]

\paragraph{Curvature-Induced Deformation.}
\[
(O2 \circ O1)_{\kappa}(x)
= O2(O1(x)) + \kappa(x).
\]

This deformation is essential for holonomy correctness.

%------------------------------------------------------------
\section{Thermodynamic Layer}

\subsection{Free-Energy Functional}

The free-energy functional is defined as
\[
F(s) = A(s) + C(s) + S(s),
\]
where $A$ is accuracy, $C$ is complexity, and $S$ is surprise.

%------------------------------------------------------------
\subsection{Accuracy Term}

Accuracy measures alignment between predicted and observed states:
\[
A(s) = \|x - \hat{x}(s)\|^2.
\]

%------------------------------------------------------------
\subsection{Complexity Term}

Complexity penalizes overly complex internal states:
\[
C(s) = \frac{\lambda}{2} \|s\|^2.
\]

Alternative sparsity-based complexity:
\[
C{\text{sparse}}(s) = \lambda \|s\|1.
\]

%------------------------------------------------------------
\subsection{Surprise Term}

Surprise measures deviation from expected sensory input:
\[
S(s) = -\log p(x \mid s).
\]

For Gaussian sensory models:
\[
S(s) =
\frac{1}{2} (x - \mus)^\top \Sigmas^{-1} (x - \mu_s)
+ \frac{1}{2} \log |\Sigma_s|.
\]

%------------------------------------------------------------
\subsection{Free-Energy Gradient}

The gradient of free energy is
\[
\nabla F(s) =
\nabla A(s) +
\nabla C(s) +
\nabla S(s).
\]

This gradient drives active inference.

%------------------------------------------------------------
\subsection{Attractor Basins}

Attractor basins are defined by
\[
\mathcal{B}(s) = \{ x \mid F(x) < F(s) \}.
\]

These basins determine the system’s thermodynamic tendencies.

%------------------------------------------------------------
\subsection{Dissipation Curves}

Dissipation between successive states is
\[
D(st, s{t+1}) = F(st) - F(s{t+1}),
\quad D(st, s{t+1}) \ge 0.
\]

This ensures thermodynamic stability.

%------------------------------------------------------------
\subsection{Active Inference Update Rule}

The full update rule is
\[
s_{t+1} =
\mathrm{ResidualMix}\Big(
    s_t,\;
    \mathcal{H}\big(
        st - \eta \cdot \mathrm{SelectedDrive}(st)
    \big)
\Big).
\]

%------------------------------------------------------------
\subsection{Jacobian of the Update Rule}

Let
\[
U(s_t) =
\mathrm{ResidualMix}\Big(
    s_t,\;
    \mathcal{H}\big(
        st - \eta \cdot \mathrm{SelectedDrive}(st)
    \big)
\Big).
\]

The Jacobian is
\[
J = \frac{\partial U}{\partial s_t}.
\]

Stability requires
\[
\|J\| \le 1.
\]

This prevents curvature blowout and holonomy distortion.
`

---

